\section{AMC 8 1985}
        \begin{enumerate}
        	\item (\textit{AMC 8 1985}) $\frac{3\times 5}{9\times 11}\times \frac{7\times 9\times 11}{3\times 5\times 7}=$


 $\textbf{(A)}\ 1 \qquad \textbf{(B)}\ 0 \qquad \textbf{(C)}\ 49 \qquad \textbf{(D)} \frac{1}{49}\ \qquad \textbf{(E)}\ 50$


 
 \textbf{Answer: }A



\textbf{Solution 1}

By the commutative property: 
 \[\frac{3\times 5}{9\times 11}\times \frac{7\times 9\times 11}{3\times 5\times 7}=\frac{3\times 5\times 7 \times 9 \times 11}{9\times 11\times 3 \times 5 \times 7}= \frac{3\times 5 \times 7 \times 9 \times 11}{3\times 5 \times 7 \times 9 \times 11}=\boxed{\textbf{(A)}~1}\]

 



\textbf{Solution 2}

Notice that the $9 \times 11$ in the denominator of the first fraction cancels with the same term in the second fraction, the $7$s in the numerator and denominator of the second fraction cancel, and the $3 \times 5$ in the numerator of the first fraction cancels with the same term in the denominator second fraction. Then everything in the expression cancels, leaving us with $\boxed{\textbf{(A)}~1}$.

 



\textbf{Solution 3 (Brute Force)}

(Note: This method is highly time consuming and should only be used as a last resort in math, also please do not do this.)$3 \times 5 \times 7 \times 9 \times 11 = 10395$

 $9 \times 11 \times 3 \times 5 \times 7 = 10395$

 Thus, the answer is $1$, or $\boxed{\textbf{(A)}\ 1}$

 



\textbf{Solution 4 (Easy, similar to solution 1)}

Notice the two fractions are multiplied, so we can consider them as 1 fraction. All terms cancel out, leaving us with $\boxed{\textbf{(A)}~1}$.~\href{https://artofproblemsolving.com/wiki/index.php?title=Shunyipanda&action=edit&redlink=1}{shunyipanda}

 



\textbf{Solution 5 (Capital Pi)}

This fraction can be simplified to 
 \[\frac{\prod_{k=1}^{5}(2k+1)}{\prod_{k=1}^{5}(2k+1)}\] The numerator and denominator are the same, so the answer is $\boxed{\textbf{(A)}\ 1}$. ~\href{https://artofproblemsolving.com/wiki/index.php?title=Shunyipanda&action=edit&redlink=1}{shunyipanda} Note: please can somebody make this look better

 


	\item (\textit{AMC 8 1985}) $90+91+92+93+94+95+96+97+98+99=$


 

$\textbf{(A)}\ 845 \qquad \textbf{(B)}\ 945 \qquad \textbf{(C)}\ 1005 \qquad \textbf{(D)}\ 1025 \qquad \textbf{(E)}\ 1045$


 
 \textbf{Answer: }B



\textbf{Solution 1}

Rearrange the terms to form $5$ pairs, each of which sums to $189$:
 \begin{align*} 90+91+\cdots +98+99 &= (90+99)+(91+98)+(92+97)+(93+96)+(94+95) \\ &= 189+189+189+189+189 \\ &= 189 \cdot 5 \\ &=\boxed{\textbf{(B)}\ 945} \end{align*}


 



\textbf{Solution 2}

Rewrite this sum as: 
 \begin{align*} 90 + 91 + 92 + \cdots + 98 + 99&= 90 + (90 + 1) + (90 + 2) + \cdots + (90 + 8) + (90 + 9) \\ &= 90 \cdot 10 + (1 + 2 + ... + 8 + 9) \\ &= 900 + \frac{9\cdot 10}{2} \\ &= \boxed{\textbf{(B)}\ 945}\end{align*}
~\href{https://artofproblemsolving.com/wiki/index.php?title=Shunyipanda&action=edit&redlink=1}{shunyipanda} (Minor edits)

 



\textbf{Solution 3}

Rewrite this sum as: 
 \begin{align*} 90 + 91 + \cdots + 98 + 99&= (100 - 10) + (100 - 9) + \cdots + (100 - 2) + (100 - 1) \\ &= 100 \cdot 10 - (10 + 9 + ... + 2 + 1) \\ &= 1000 - \frac{10 \cdot 11}{2} \\ &= \boxed{\textbf{(B)}\ 945}\end{align*}


 


 



\textbf{Solution 4 (Arithmetic Series Formula)}

Recall that the sum of an arithmetic series is $\frac{n}{2} (a_1+a_n)$, where $n$ is the number of terms, $a_1$ is the first term and $a_n$ is the last term.

 Therefore, the sum of this series is $\frac{10}{2}(90+99)=\boxed{\textbf{(B)}\ 945}$.

 



\textbf{Solution 5 (Difference of Triangular Numbers)}

Notice that this expression is the sum of the first $99$ positive integers minus the sum of the first $89$ positive integers, that is,  $\frac{99\cdot100}{2} - \frac{89\cdot90}{2} = 4950 - 4005 = \boxed{\textbf{(B)}\ 945}$.

 



\textbf{Solution 6 (Answer Choices, Estimation)}

Since all the terms are between $90$ and $100$, the sum must be between $10 \cdot 90 = 900$ and $10 \cdot 100 = 1000$. Only  $\boxed{\textbf{(B)}\ 945}$ is between $900$ and $1000$.

 



\textbf{Solution 7 (Brute Force)}

Adding all the terms yields $\boxed{\textbf{(B)}\ 945}$.

 


	\item (\textit{AMC 8 1985}) $\frac{10^7}{5\times 10^4}=$


 

$\textbf{(A)}\ .002 \qquad \textbf{(B)}\ .2 \qquad \textbf{(C)}\ 20 \qquad \textbf{(D)}\ 200 \qquad \textbf{(E)}\ 2000$


 
 \textbf{Answer: }D



\textbf{Solution 1}

Cancelling powers of $10$ yields:
 \[\frac{10^7}{5 \times 10^4} = \frac{10^3}{5} = \frac{1000}{5} = \boxed{\textbf{(D) } 200}\] (Reverted a Mistake)

 



\textbf{Solution 2 (Brute Force)}

$\frac{10^7}{5\times 10^4}=\frac{10000000}{50000}=\boxed{\textbf{(D)}\ 200}$

 



\textbf{Solution 3 (Reasoning)}

As we can see, the numerator is way larger than the denominator so answer choices $\textbf{(A)}\ .002$, $\textbf{(B)}\ .2$, and $\textbf{(C)}\ 20$ are eliminated. $\frac{10^7}{10^4}$ is $1000$, but in this question the expression, $\frac{10^7}{5 \times 10^4}$, is smaller than $1000$ since the denominator is larger than $10^4$. Therefore the answer is $\boxed{\textbf{(D)}\ 200}$.~\href{https://artofproblemsolving.com/wiki/index.php?title=Shunyipanda&action=edit&redlink=1}{shunyipanda}

 


	\item (\textit{AMC 8 1985}) The area of polygon $ABCDEF$, in square units, is


 $\textbf{(A)}\ 24 \qquad \textbf{(B)}\ 30 \qquad \textbf{(C)}\ 46 \qquad \textbf{(D)}\ 66 \qquad \textbf{(E)}\ 74$

                    \begin{figure}[H]
                        \centering
                        \includegraphics[width=0.5\textwidth]{Problem/AMC_8_AJHSME_1985_4_Problem_Figure_1.png}
                    \end{figure}
                    


 
 \textbf{Answer: }C



\textbf{Solution 1 (Adding Rectangular Areas)}


                    \begin{figure}[H]
                        \centering
                        \includegraphics[width=0.5\textwidth]{Solution/AMC_8_AJHSME_1985_4_Solution_Figure_1.png}
                    \end{figure}
                    

 Extend $\overline{FE}$ to meet $\overline{BC}$ at $G$, dividing $ABCDEF$ into rectangles $ABGF$ and $EGCD$.  Then $BG = AF = 5$ and $GC = 9 - 5 = 4$.  

 So $ABGF$ has area $6 \cdot 5 = 30$, while $EGCD$ has area $4 \cdot 4 = 16$.  Therefore, the total area of $ABCDEF$ is $30 + 16 = \boxed{\textbf{(C)}\ 46}$.

 



\textbf{Solution 2 (Subtracting Rectangular Areas)}


                    \begin{figure}[H]
                        \centering
                        \includegraphics[width=0.5\textwidth]{Solution/AMC_8_AJHSME_1985_4_Solution_Figure_2.png}
                    \end{figure}
                    

 Extend $\overline{CD}$ and $\overline{AF}$ to meet at $G$.  Then $FG = 9 - 5 = 4$ and $GD = 6 - 4 = 2$.  

 The area of $ABCDEF$ is the difference between the area of rectangle $ABCG$ and the area of rectangle $FEDG$.

 Since the area of $ABCG$ is $6 \cdot 9 = 54$, while the area of $FEDG$ is $4 \cdot 2 = 8$, the area of $ABCDEF$ is $54 - 8 = \boxed{\textbf{(C)}\ 46}$.

 



\textbf{Solution 3 (Trapezoids)}


                    \begin{figure}[H]
                        \centering
                        \includegraphics[width=0.5\textwidth]{Solution/AMC_8_AJHSME_1985_4_Solution_Figure_3.png}
                    \end{figure}
                    

 First, notice $FE = 6 - 4 = 2$ and $ED = 9 - 5 = 4$.  

 Draw $\overline{BE}$ to divide $ABCDEF$ into trapezoids $ABEF$ and $BCDE$.  Then the area of $ABEF = \frac{5(6+2)}{2} = 20$ and the area of $BCDE$ is $\frac{4(4+9)}{2} = 26$, so the area of $ABCDEF$ is $20+26=\boxed{\textbf{(C)}\ 46}$.

 


	\item (\textit{AMC 8 1985}) The bar graph shows the grades in a mathematics class for the last grading period. If $A$, $B$, $C$, and $D$ are satisfactory grades, what fraction of the grades shown in the graph are satisfactory?


 $\textbf{(A)}\ \frac{1}{2} \qquad \textbf{(B)}\ \frac{2}{3} \qquad \textbf{(C)}\ \frac{3}{4} \qquad \textbf{(D)}\ \frac{4}{5} \qquad \textbf{(E)}\ \frac{9}{10}$


 


                    \begin{figure}[H]
                        \centering
                        \includegraphics[width=0.5\textwidth]{Problem/AMC_8_AJHSME_1985_5_Problem_Figure_1.png}
                    \end{figure}
                    


 
 \textbf{Answer: }C



\textbf{Solution 1}

Of $5 + 4 + 3 + 3 + 5 = 20$ students in the class, $5 + 4 + 3 + 3 = 15$ received satisfactory grades.  So $\frac{15}{20}=\boxed{\textbf{(C)}\ \frac{3}{4}}$ of the class received satisfactory grades.

 



\textbf{Solution 2}

Of $5 + 4 + 3 + 3 + 5 = 20$ students in the class, $5$ did not receive satisfactory grades.  That is, $\frac{5}{20} = \frac{1}{4}$ of the students did not receive satisfactory grades, and thus $1 - \frac{1}{4} = \boxed{\textbf{(C)}\ \frac{3}{4}}$ did.

 


	\item (\textit{AMC 8 1985}) A ream of paper containing $500$ sheets is $5$ cm thick. Approximately how many sheets of this type of paper would there be in a stack $7.5$ cm high?


 $\textbf{(A)}\ 250 \qquad \textbf{(B)}\ 550 \qquad \textbf{(C)}\ 667 \qquad \textbf{(D)}\ 750 \qquad \textbf{(E)}\ 1250$


 
 \textbf{Answer: }D



\textbf{Solution 1}

Each sheet is $\frac{5}{500} = 0.01\ \text{cm}$ thick, so a $7.5\ \text{cm}$ stack contains $\frac{7.5}{0.01} = \boxed{\textbf{(D)}\ 750}$ sheets.

 



\textbf{Solution 2}

Since each sheet has the same thickness, the number of sheets in a stack is directly proportional to the thickness of a stack.  

 Let $x$ be the number of sheets in a $7.5\ \text{cm}$ stack.Then $\frac{5}{500}=\frac{7.5}{x}$, so $5x=3750$ and thus $x = \boxed{\textbf{(D)}\ 750}$.

 



\textbf{Solution 3}

Using dimensional analysis with the conversion factor $\frac{500\ \text{sheets}}{5\ \text{cm}}$ yields $(7.5\ \text{cm})\left(\frac{500\ \text{sheets}}{5\ \text{cm}}\right) = \boxed{\textbf{(D)}\ 750}\ \text{sheets}$.

 



\textbf{Solution 4}

Each centimeter of a stack contains $\frac{500}{5} = 100$ sheets of paper, so a $7.5\ \text{cm}$ stack must contain $7.5 \cdot 100 = \boxed{\textbf{(D)}\ 750}$ sheets of paper.~\href{https://artofproblemsolving.com/wiki/index.php?title=Shunyipanda&action=edit&redlink=1}{shunyipanda}

 



\textbf{Solution 5}

Since a $7.5\ \text{cm}$ stack is $1.5$ times as thick as a $5\ \text{cm}$ stack, a $7.5\ \text{cm}$ stack contains $1.5 \cdot 500  = \boxed{\textbf{(D)}\ 750}$ sheets of paper. ~\href{https://artofproblemsolving.com/wiki/index.php?title=Shriya2}{Shriya2}

 


	\item (\textit{AMC 8 1985}) A “stair‑step” figure is made of alternating black and white squares in each row. Rows $1$ through $4$ are shown. All rows begin and end with a white square. The number of black squares in the $37$th row is


 $\textbf{(A)}\ 34 \qquad \textbf{(B)}\ 35 \qquad \textbf{(C)}\ 36 \qquad \textbf{(D)}\ 37 \qquad \textbf{(E)}\ 38$


 
 \textbf{Answer: }C



\textbf{Solution 1}

The first row has $0$ black squares, while each row after the first has $1$ more black square than the previous row.  Therefore, the second row has $1$ black square, the third row has $2$ black squares, and so on, until the $37^{\text{th}}$ row has $\boxed{\textbf{(C)}\ 36}$ black squares.

 



\textbf{Solution 2}

Let $f(n)$ be the number of black squares in the $n^{\text{th}}$ row.  Since each row has $1$ more black square than the previous row, $f(n)$ is linear with slope $1$, that is, $f(n) = n + b$ for some constant $b$.  Since $f(1) = 0$, $0 = 1 + b$ and $b = -1$.  So $f(n) = n-1$ and the $37^{\text{th}}$ row has $f(37) = 37 - 1 = \boxed{\textbf{(C)}\ 36}$ black squares.

 


	\item (\textit{AMC 8 1985}) If $a = - 2$, the largest number in the set $\{- 3a, 4a, \frac {24}{a}, a^2, 1 \}$ is


 $\textbf{(A)}\ -3a \qquad \textbf{(B)}\ 4a \qquad \textbf{(C)}\ \frac {24}{a} \qquad \textbf{(D)}\ a^2 \qquad \textbf{(E)}\ 1$


 
 \textbf{Answer: }A



\textbf{Solution 1}

Substituting $a = -2$ yields $\{ - 3a, 4a, \frac {24}{a}, a^2, 1\} = \{ (-3)(-2), 4(-2), \frac{24}{-2}, (-2)^2, 1 \}= \{ 6, -8, -12, 4, 1 \}$.

 The largest element of this set is therefore $6$, that is, $\boxed{\textbf{(A)} -3a}$.

 



\textbf{Solution 2}

Since $a$ is negative, $-3a$, $a^2$, and $1$ must be positive while $4a$ and $\frac{24}{a}$ are negative.  Evaluating the positive expressions shows that $\boxed{\textbf{(A)} -3a} = 6$ is the largest.

 


	\item (\textit{AMC 8 1985}) The product of the 9 factors $\Big(1 - \frac12\Big)\Big(1 - \frac13\Big)\Big(1 - \frac14\Big)\cdots\Big(1 - \frac {1}{10}\Big) =$


 $\textbf{(A)}\ \frac {1}{10} \qquad \textbf{(B)}\ \frac {1}{9} \qquad \textbf{(C)}\ \frac {1}{2} \qquad \textbf{(D)}\ \frac {10}{11} \qquad \textbf{(E)}\ \frac {11}{2}$


 
 \textbf{Answer: }A



\textbf{Solution 1}

Simplifying each factor and cancelling yields $\Big(1 - \frac12\Big)\Big(1 - \frac13\Big)\Big(1 - \frac14\Big)\cdots\Big(1 - \frac {1}{10}\Big) =\left(\frac{1}{2}\right)\left(\frac{2}{3}\right)\left(\frac{3}{4}\right)\cdots\left(\frac{9}{10}\right) = \frac{1\cdot2\cdot3\cdot4\cdot5\cdot6\cdot7\cdot8\cdot9}{2\cdot3\cdot4\cdot5\cdot6\cdot7\cdot8\cdot9\cdot10} = \boxed{\textbf{(A)}\ \frac{1}{10}}$.

 



\textbf{Solution 2 (Induction, Generalization)}

It can be proven by induction that $\left(1 - \frac{1}{2}\right)\cdots\left(1 - \frac{1}{n}\right) = \frac{1}{n}$ for all integers $n \geq 2$.

 For the base case, if $n = 2$, then this product is simply $1 - \frac{1}{2} = \frac{1}{2}$.

 For the inductive case, assume that $\left(1 - \frac{1}{2}\right)\cdots\left(1-\frac{1}{k}\right) = \frac{1}{k}$.  Then $\left(1 - \frac{1}{2}\right)\cdots\left(1-\frac{1}{k+1}\right) = \left(1 - \frac{1}{2}\right)\cdots\left(1-\frac{1}{k}\right)\left(1-\frac{1}{k+1}\right) = \left(\frac{1}{k}\right)\left(\frac{k}{k+1}\right) = \frac{k}{k(k+1)} = \frac{1}{k+1}$.  This completes the induction.

 Applying this lemma for $n = 10$ yields $\boxed{\textbf{(A)}\ \frac{1}{10}}$.

 


	\item (\textit{AMC 8 1985}) The fraction halfway between $\frac{1}{5}$ and $\frac{1}{3}$ (on the number line) is


 
                    \begin{figure}[H]
                        \centering
                        \includegraphics[width=0.5\textwidth]{Problem/AMC_8_AJHSME_1985_10_Problem_Figure_1.png}
                    \end{figure}
                    


 $\textbf{(A)}\ \frac{1}{4} \qquad \textbf{(B)}\ \frac{2}{15} \qquad \textbf{(C)}\ \frac{4}{15} \qquad \textbf{(D)}\ \frac{53}{200} \qquad \textbf{(E)}\ \frac{8}{15}$


 
 \textbf{Answer: }C



\textbf{Solution 1}

The fraction halfway between $\frac{1}{5}$ and $\frac{1}{3}$ is their average, which is $\frac{1}{2}\left(\frac{1}{5}+\frac{1}{3}\right) = \frac{1}{2}\left(\frac{8}{15}\right) = \boxed{\textbf{(C)}\ \frac{4}{15}}$

 



\textbf{Solution 2}


                \begin{itemize}
                    \item For any two fractions that have the property $\frac{a}{b}$ and $\frac{a}{c}$ where $a$ is equal for both fractions;
                \end{itemize}


                Set aside the a. Take the average of b and c. For this example, the average of 3 and 5 is 4. Then multiply those numbers. 3*5=15. The number that is exactly in between $\frac{a}{b}$ and $\frac{a}{c}$ will be $\frac{a*((b+c)/2)}{b*c}$.

 In this example, it would be $\frac{1((3+5)/2)}{3*5}$, which simplifies to $\frac{4}{15}$.

 Therefore, the answer is $\boxed{\text{C}}$

 Edit by mathmagical~

 


	\item (\textit{AMC 8 1985}) A piece of paper containing six joined squares labeled as shown in the diagram is folded along the edges of the squares to form a cube.  The label of the face opposite the face labeled $\text{X}$ is 


 
                    \begin{figure}[H]
                        \centering
                        \includegraphics[width=0.5\textwidth]{Problem/AMC_8_AJHSME_1985_11_Problem_Figure_1.png}
                    \end{figure}
                    


 $\textbf{(A)}\ \text{Z} \qquad \textbf{(B)}\ \text{U} \qquad \textbf{(C)}\ \text{V} \qquad \textbf{(D)}\ \ \text{W} \qquad \textbf{(E)}\ \text{Y}$


 
 \textbf{Answer: }E



\textbf{Solution 1}

Notice that two faces of a cube are opposite each other if and only if they share no vertices.  The diagram shows that $\text{X}$ shares at least one vertex with each of $\text{U}$, $\text{V}$, $\text{W}$, and $\text{Z}$, so by process of elimination $\text{X}$ must be opposite $\boxed{\textbf{(E)}\  \text{Y}}$.

 



\textbf{Solution 2}

To find the face opposite $\text{X}$, we can find the faces sharing an edge with $\text{X}$, so the only face remaining will be the opposite face.

 Clearly, $\text{V}$ and $\text{Z}$ share an edge with $\text{X}$.  Also, the faces $\text{V}$, $\text{X}$, and $\text{W}$ share a common vertex, therefore $\text{X}$ shares an edge with $\text{W}$.  Similarly, the faces $\text{U}$, $\text{V}$, and $\text{X}$ share a common vertex, so $\text{X}$ shares an edge with $\text{U}$.  

 The only face $\text{X}$ doesn't share an edge with is $\text{Y}$, which is choice $\boxed{\text{E}}$

 


	\item (\textit{AMC 8 1985}) A square and a triangle have equal perimeters.  The lengths of the three sides of the triangle are $6.2 \text{ cm}$, $8.3 \text{ cm}$ and $9.5 \text{ cm}$.  The area of the square is


 $\textbf{(A)}\ 24\text{ cm}^2 \qquad \textbf{(B)}\ 36\text{ cm}^2 \qquad \textbf{(C)}\ 48\text{ cm}^2 \qquad \textbf{(D)}\ 64\text{ cm}^2 \qquad \textbf{(E)}\ 144\text{ cm}^2$


 
 \textbf{Answer: }B



\textbf{Solution}

The perimeter of the triangle is $6.2+8.3+9.5=24\text{ cm}$.  Therefore, the perimeter of the square is also $24 \text{ cm}$, so its side length is $\frac{24}{4}=6 \text{ cm}$ and its area is $6^2= \boxed{\textbf{(B)}\ 36 \text{ cm}^2}$.

 


	\item (\textit{AMC 8 1985}) If you walk for $45$ minutes at a rate of $4 \text{ mph}$ and then run for $30$ minutes at a rate of $10\text{ mph,}$ how many miles will you have gone at the end of one hour and $15$ minutes?


 $\textbf{(A)}\ 3.5\text{ miles} \qquad \textbf{(B)}\ 8\text{ miles} \qquad \textbf{(C)}\ 9\text{ miles} \qquad \textbf{(D)}\ 25\frac{1}{3}\text{ miles} \qquad \textbf{(E)}\ 480\text{ miles}$


 
 \textbf{Answer: }B



\textbf{Solution}

Since you walked for $45$ minutes, or $\frac{3}{4}$ hour, you walked $\frac{3}{4}\cdot 4 = 3$ miles.  Since you ran for $30$ minutes, or $\frac{1}{2}$ hour, you ran $\frac{1}{2}\cdot 10 = 5$ miles.  Therefore, you traveled a total of $3+5=\boxed{\textbf{(B)}\ 8\text{ miles}}$.

 


	\item (\textit{AMC 8 1985}) The difference between a $6.5\%$ sales tax and a $6\%$ sales tax on an item priced at $\$20$ before tax is


 $\text{(A)}$ $\$.01$


 $\text{(B)}$ $\$.10$


 $\text{(C)}$ $\$ .50$


 $\text{(D)}$ $\$ 1$


 $\text{(E)}$ $\$10$


 
 \textbf{Answer: }B



\textbf{Solution 1}

With a $6.5\%$ sales tax, an item costs $1.065$ times its listed price, while with a $6\%$ sales tax, it costs $1.06$ times its listed price.  The difference between these prices is therefore $1.065 - 1.06 = .005$ times its listed price.  For this item, this difference is $.005 \cdot \$20 = \boxed{\textbf{(B)}\ \$.10}$.

 



\textbf{Solution 2}

With a $6.5\%$ sales tax, the item costs $1.065 \cdot \$20 = \$21.30$.  With a $6\%$ sales tax, the item costs $1.06 \cdot \$20 = \$21.20$.  So the difference is $\$21.30 - \$21.20 = \boxed{\textbf{(B)}\ \$.10}$.

 


	\item (\textit{AMC 8 1985}) How many whole numbers between $100$ and $400$ contain the digit $2$?


 $\textbf{(A)}\ 100 \qquad \textbf{(B)}\ 120 \qquad \textbf{(C)}\ 138 \qquad \textbf{(D)}\ 140 \qquad \textbf{(E)}\ 148$


 



 
 \textbf{Answer: }C



\textbf{Solution 1 (Complementary Counting)}

Since $400$ does not contain the digit $2$, only the set of integers from $100$ to $399$ needs to be considered.  Notice that this is the set of all three-digit integers with $1$, $2$, or $3$ as their first digit.  

 Since there are $3$ choices for the first digit and $10$ choices for each of the second and third digits, this set contains $3 \cdot 100 \cdot 100 = 300$ numbers.

 For a three-digit integer between $100$ and $399$ that does not contain the digit $2$, there are $3-1=2$ choices for the first digit and $10-1=9$ choices for each of the second and third digit, for a total of $2 \cdot 9 \cdot 9 = 162$ numbers.

 Therefore, $300 - 162 = \boxed{\textbf{(C)}\ 138}$ of these numbers contain the digit $2$.

 



\textbf{Solution 2}

All $100$ integers from $200$ to $299$ start with $2$.  

 Of the integers from $100$ to $400$ not starting with $2$, $20$ have a second digit of $2$: $\{120, 121, 122, \cdots, 129, 320, 321, 322, \cdots, 329\}$, while $20$ have a third digit of $2$: $\{102, 112, 122, \cdots, 192, 302, 312, 322, \cdots, 392\}$. Of these, $122$ and $322$ have both a second and third digit of $2$, so a total of $20 + 20 - 2 = 38$ integers between $100$ and $400$ not starting with $2$ contain the digit $2$.

 Therefore, a total of $100 + 38 = \boxed{\textbf{(C)}\ 138}$ integers between $100$ and $400$ contain the digit $2$.

 


	\item (\textit{AMC 8 1985}) The ratio of boys to girls in Mr. Brown's math class is $2:3$.  If there are $30$ students in the class, how many more girls than boys are in the class?


 $\text{(A)}\ 1 \qquad \text{(B)}\ 3 \qquad \text{(C)}\ 5 \qquad \text{(D)}\ 6 \qquad \text{(E)}\ 10$


 
 \textbf{Answer: }D



\textbf{Solution 1}

Let $x$ be half the number of boys.  Then there are $2x$ boys and $3x$ girls. So $2x+3x=5x=30$ and thus $x = 6$. Therefore, between the number of girls and the number of boys is $3x-2x=x=\boxed{\textbf{(D)}\ 6}$.

 



\textbf{Solution 2}

Since the ratio of boys to girls is $2 : 3$, $\frac{2}{2+3} = \frac{2}{5}$ of the students are boys and $\frac{3}{2+3} = \frac{3}{5}$ are girls.  Thus the difference between the number of boys and the number of girls is $\frac{3}{5} - \frac{2}{5} = \frac{1}{5}$ of the number of students in the class, that is $\frac{1}{5} \cdot 30 = \boxed{\textbf{(D)}\ 6}$.

 


	\item (\textit{AMC 8 1985}) If your average score on your first six mathematics tests was $84$ and your average score on your first seven mathematics tests was $85$, then your score on the seventh test was


 $\textbf{(A)}\ 86 \qquad \textbf{(B)}\ 88 \qquad \textbf{(C)}\ 90 \qquad \textbf{(D)}\ 91 \qquad \textbf{(E)}\ 92$


 
 \textbf{Answer: }D



\textbf{Solution 1}

Since your average score on your first $6$ tests is $84$, your total score on your first $6$ tests is $6\cdot84=504$. 

 Since your average score on your first $7$ tests is $85$, your total score on your first $7$ tests is $7\cdot85= 595$.  

 Therefore, your score on your $7^{\text{th}}$ test was $595-504=\boxed{\textbf{(D)}\ 91}$.

 



\textbf{Solution 2}

Let's remove the condition that the average of the first seven tests is $85$, and say the 7th test score was a $0$. Then, the average of the first seven tests would be 
 \[\frac{6\times 84}{7}=72\]

 If we increase the seventh test score by $7$, the average will increase by $\frac{7}{7}=1$. We need the average to increase by $85-72=13$, so the seventh test score is $7\times 13=91$ more than $0$, which is clearly $\boxed{\textbf{(D)}\ 91}$.

 



\textbf{Solution 3 (Answer Choices)}

Since $84$ is divisible by $7$, your total score on your first $6$ tests, that is, $6 \cdot 84$, is divisible by $7$.  Also, your total score on your first $7$ tests, that is, $85 \cdot 7$, must be divisible by $7$.  Therefore, your score on your $7^{\text{th}}$ test is also divisible by $7$. Of the answer choices, only $\boxed{\textbf{(D)}\ 91}$ is divisible by $7$.

 


	\item (\textit{AMC 8 1985}) Nine copies of a certain pamphlet cost less than $\$10.00$ while ten copies of the same pamphlet (at the same price) cost more than $\$11.00$.  How much does one copy of this pamphlet cost?


 $\textbf{(A)}$ $\$1.07$


 $\textbf{(B)}$ $\$1.08$


 $\textbf{(C)}$ $\$1.09$


 $\textbf{(D)}$ $\$1.10$


 $\textbf{(E)}$ $\$1.11$


 
 \textbf{Answer: }E



\textbf{Solution}

Let $p$ be the cost of one pamphlet in dollars. The first restriction implies $9p<10$, that is, $p<1.1111...$, while the section restriction implies $10p>11$, that is, $p>1.10$.  Thus $1.10 < p < 1.1111...$, and thus $p = \boxed{\textbf{(E)}\ \$1.11}$, since the price must be a whole number of cents.

 Note: the second restriction alone eliminates all the incorrect answer choices.

 


	\item (\textit{AMC 8 1985}) If the length and width of a rectangle are each increased by $10\%$, then the perimeter of the rectangle is increased by


 $\textbf{(A)}\ 1\% \qquad \textbf{(B)}\ 10\% \qquad \textbf{(C)}\ 20\% \qquad \textbf{(D)}\ 21\% \qquad \textbf{(E)}\ 40\%$


 
 \textbf{Answer: }B



\textbf{Solution 1}

Let $w$, $l$, and $p$ be the original width, length, and perimeter of the rectangle, respectively. Then $p = 2(w+l)$. 

 The width and length are increased to $1.1w$ and $1.1l$, respectively. Therefore, the new perimeter is $2(1.1w+1.1l)=2\cdot1.1(w+l) = 1.1 \cdot 2(w+l) = 1.1p$, or $\boxed{\textbf{(B)}\ 10\%}$ more than the original perimeter.

 



\textbf{Solution 2}

Since each side is scaled by a factor of $1.1$, the new rectangle is similar to the original, and therefore its perimeter is also scaled by a factor of $1.1$, that is, a $\boxed{\textbf{(B)}\ 10\%}$ increase.

 



\textbf{Solution 3}

Suppose the original rectangle is a square of side length $10$, so its perimeter is $4 \cdot 10 = 40$.  After scaling, its sides will have length $1.1 \cdot 10 = 11$, so its perimeter will be $11 \cdot 4 = 44$.  Therefore, the new perimeter is $\frac{44}{40} = 1.1$ times the original perimeter, a $\boxed{\textbf{(B)}\ 10\%}$ increase.

 


	\item (\textit{AMC 8 1985}) In a certain year, January had exactly four Tuesdays and four Saturdays.  On what day did January $1$ fall that year?


 $\textbf{(A)}\ \text{Monday} \qquad \textbf{(B)}\ \text{Tuesday} \qquad \textbf{(C)}\ \text{Wednesday} \qquad \textbf{(D)}\ \text{Friday} \qquad \textbf{(E)}\ \text{Saturday}$


 
 \textbf{Answer: }C



\textbf{Solution 1}

January has $31$ days, that is, $4$ weeks followed by $3$ extra days.  Since there are only $4$ Tuesdays and $4$ Saturdays, the $3$ extra days must include neither Tuesday nor Saturday.  This is only possible if the extra days are Wednesday through Friday.  Therefore, January 29 (and thus January 1) is a $\boxed{\textbf{(C)}\ \text{Wednesday}}$.

 



\textbf{Solution 2 (Answer Choices)}

If January $1$ was on Monday or Tuesday, then there would be $5$ Tuesdays.  If January $1$ was on Friday or Saturday, there would be $5$ Saturdays.  This eliminates all choices except $\boxed{\textbf{(C)}\ \text{Wednesday}}$.

 


	\item (\textit{AMC 8 1985}) Mr. Green receives a $10\%$ raise every year.  His salary after four such raises has gone up by what percent?


 $\textbf{(A)}\ \text{less than }40\% \qquad \textbf{(B)}\ 40\% \qquad \textbf{(C)}\ 44\% \qquad \textbf{(D)}\ 45\% \qquad \textbf{(E)}\ \text{more than }45\%$


 
 \textbf{Answer: }E



\textbf{Solution 1}

Each year, Mr. Green's salary is multiplied by $1.1$.  Therefore, after $4$ years, his salary will be $1.1^4 = 1.4641$ times his current salary, for a total increase of $46.41\%$, which is $\boxed{\textbf{(E)}\ \text{more than }45\% }$.

 



\textbf{Solution 2}

Suppose Mr. Green's salary is originally $\$100.00$. After $1$ year, his salary will be $\$110.00$. After $2$ years, his salary will be $\$121.00$.  After $3$ years, his salary will be $\$133.10$.  After $4$ years, his salary will be $\$146.41$, which is $1.4641$ times his original salary.  This is a $46.41\%$ increase, so the answer is $\boxed{\textbf{(E)}\ \text{more than 45\%}}$. 

 


	\item (\textit{AMC 8 1985}) Assume every 7-digit whole number is a possible telephone number except those that begin with $0$ or $1$.  What fraction of telephone numbers begin with $9$ and end with $0$?


 $\textbf{(A)}\ \frac{1}{63} \qquad \textbf{(B)}\ \frac{1}{80} \qquad \textbf{(C)}\ \frac{1}{81} \qquad \textbf{(D)}\ \frac{1}{90} \qquad \textbf{(E)}\ \frac{1}{100}$


 
 \textbf{Answer: }B



\textbf{Solution 1}

For a valid phone number, there are $10-2=8$ possibilities for the first digit and $10$ for each of the other $6$ digits.  Therefore, there are $8 \cdot 10^6 = 8000000$ valid phone numbers.

 For a phone number that starts with $9$ and ends with $0$, there are $10$ possibilities for each of the middle $5$ digits.  Therefore, there are $10^5 = 100000$ phone numbers that start with $9$ and end with $0$.

 Therefore, $\frac{100000}{8000000} = \boxed{\textbf{(B)}\ \frac{1}{80}}$ of valid phone numbers start with $9$ and end with $0$.

 



\textbf{Solution 2}

Notice that selecting a random valid phone number is identical to randomly and independently selecting each digit of the phone number.

 There are $10-2=8$ possibilities for the first digit and $10$ possibilities for the last digit, so the probability of selecting a phone number with first digit $9$ and last digit $0$ is $\frac{1}{8} \cdot \frac{1}{10} = \boxed{\textbf{(B)}\ \frac{1}{80}}$.

 


	\item (\textit{AMC 8 1985}) King Middle School has $1200$ students.  Each student takes $5$ classes a day.  Each teacher teaches $4$ classes.  Each class has $30$ students and $1$ teacher.  How many teachers are there at King Middle School?  


 $\textbf{(A)}\ 30 \qquad \textbf{(B)}\ 32 \qquad \textbf{(C)}\ 40 \qquad \textbf{(D)}\ 45 \qquad \textbf{(E)}\ 50$


 
 \textbf{Answer: }E



\textbf{Solution 1}

Each of the $1200$ students attends $5$ class sessions, so they collectively attend a total of $1200 \cdot 5 = 6000$ class sessions.  Since $30$ students attend each class session, there must be $\frac{6000}{30}=200$ class sessions held.  Since each class has $1$ teacher, the teachers collectively host a total of $200 \cdot 1 = 200$ class sessions.  Since each teacher hosts $4$ class sessions, there must be $\frac{200}{4} = \boxed{\textbf{(E)}\ 50}$ teachers.

 



\textbf{Solution 2 (Similar to above solution)}

Each class is $30$ students so, there are $40$ classes in total, that need to take $200$ periods every day.

 As each teacher is only taking $4$ periods, there will be needed ${200 \over 4} = 50$ teachers.

 


	\item (\textit{AMC 8 1985}) In a magic triangle, each of the six whole numbers $10\text{--}15$ is placed in one of the circles so that the sum, $S$, of the three numbers on each side of the triangle is the same.  The largest possible value for $S$ is


 
                    \begin{figure}[H]
                        \centering
                        \includegraphics[width=0.5\textwidth]{Problem/AMC_8_AJHSME_1985_24_Problem_Figure_1.png}
                    \end{figure}
                    


 $\textbf{(A)}\ 36 \qquad \textbf{(B)}\ 37 \qquad \textbf{(C)}\ 38 \qquad \textbf{(D)}\ 39 \qquad \textbf{(E)}\ 40$


 



 
 \textbf{Answer: }D



\textbf{Solution 1}

Let $x$ be the sum of the vertices and $y$ be the sum of the midpoints.  Then $x + y = 10 + 11 + 12 + 13 + 14 + 15 = 75$.  Also, since each vertex is included in $2$ sides while each midpoint is included in $1$ side, $3S = 2x + y = x + (x+y) = x + 75$, so $S = \frac{x}{3} + 25$.  Therefore, maximizing $x$ also maximizes $S$.  

 Since $x \leq 13 + 14 + 15 = 42$, $S \leq \frac{42}{3} + 25 = 39$.  Achieving $S = \boxed{\textbf{(D)}\ 39}$ is possible using the following configuration:

 
                    \begin{figure}[H]
                        \centering
                        \includegraphics[width=0.5\textwidth]{Solution/AMC_8_AJHSME_1985_24_Solution_Figure_1.png}
                    \end{figure}
                    

 



\textbf{Solution 2 (Answer Choices)}

As in Solution 1, notice that if $x$ is the sum of the vertices, then $S = \frac{x}{3} + 25$.  Trying $S = 40$, we find that $x=45$ which is too large.  Trying $S = \boxed{\textbf{(D)}\ 39}$, we find that $x = 42$ which works as shown in the configuration in Solution 1.

 -scthecool, edited by j314andrews

 


	\item (\textit{AMC 8 1985}) Five cards are lying on a table as shown.


 
 \[\begin{matrix} & \qquad & \boxed{\texttt{P}} & \qquad & \boxed{\texttt{Q}} \\  \\ \boxed{\texttt{3}} & \qquad & \boxed{\texttt{4}} & \qquad & \boxed{\texttt{6}} \end{matrix}\]


 Each card has a letter on one side and a whole number on the other side.  Jane said, "If a vowel is on one side of any card, then an even number is on the other side."  Mary showed Jane was wrong by turning over one card.  Which card did Mary turn over?


 $\textbf{(A)}\ 3 \qquad \textbf{(B)}\ 4 \qquad \textbf{(C)}\ 6 \qquad \textbf{(D)}\ \text{P} \qquad \textbf{(E)}\ \text{Q}$


 
 \textbf{Answer: }A



\textbf{Solution 1}

Recall that a counterexample to the statement "If $X$, then $Y$" is a situation in which $X$ is true but $Y$ is false.  Therefore, a counterexample to Jane's statement would be a card with a vowel on one side and an odd number on the other side. Of the cards shown, only $\boxed{\textbf{(A)}\ 3}$ could fit this description.

 



\textbf{Solution 2 (Exploitation of Multiple Choice)}

The only attributes of letters and numbers mentioned in Jane's statement are vowel/consonant and odd/even.  Only one answer is correct, so since 4 and 6 are both even numbers, and P and Q are both consonants, none of those answers could be correct.  Thus $\boxed{\textbf{(A) }3}$ must be correct.

 RCFX - Minor Changes

 


\end{enumerate}